% Generated draft: requires mathematical and source review.
\documentclass[11pt]{article}
\usepackage[margin=25mm]{geometry}
\usepackage{fontspec}
\setmainfont{Latin Modern Roman}
\usepackage{amsmath,amssymb,mathtools,unicode-math}
\setmathfont{Latin Modern Math}
\usepackage{enumitem,needspace,xurl,hyperref}
\setlength{\parindent}{0pt}
\setlength{\parskip}{5pt}
\setlength{\emergencystretch}{3em}
\newcommand{\N}{\mathbb{N}}
\newcommand{\Z}{\mathbb{Z}}
\newcommand{\Q}{\mathbb{Q}}
\newcommand{\R}{\mathbb{R}}
\newcommand{\C}{\mathbb{C}}
\newcommand{\F}{\mathbb{F}}
\newcommand{\A}{\mathbb{A}}
\newcommand{\PP}{\mathbb{P}}
\newcommand{\E}{\mathbb{E}}
\newcommand{\eps}{\varepsilon}
\newcommand{\dd}{\,\mathrm{d}}
\begin{document}
\section*{UnsolvedMath ID: 30004774}
\textbf{Source problem code:} catalogue entry
\subsection{Definitions and mathematical statement}
The supplied record attributes the following conjectural formula to Hiraga--Ichino--Ikeda--Langlands. Its local setting and normalizations are not fully specified, so the statement below retains those unresolved qualifications.

Let \(H\) be a reductive group in the intended local setting, and let \(\mathcal T_H\) denote the intended space of tempered representations. The record does not specify the base field, the precise class of groups, or whether \(\mathcal T_H\) consists of unitary equivalence classes of irreducible tempered representations. Write \(d_H(\sigma)\) for the Plancherel density appearing in the source. A density requires a reference measure and a normalization of Haar measure; neither is supplied. Likewise, the measure governing ``almost every'' is unspecified.

For \(\sigma\in\mathcal T_H\), retain the source symbols \(S_\sigma\), \(\rho_\sigma\), and \(\operatorname{Ad}_H\). The quotient \(\dim(\rho_\sigma)/|S_\sigma|\) requires a finite group \(S_\sigma\) and a finite-dimensional representation \(\rho_\sigma\). Their precise definitions, their association with \(\sigma\), and any required Langlands parameter or enhancement are absent from the record. The intended adjoint representation and the associated local gamma factor are also not defined there.

To explain the Laurent-coefficient terminology, a nonzero meromorphic factor near \(s=0\) has an expansion
\[
\gamma(s,\sigma,\operatorname{Ad}_H)
 =\sum_{j=m}^{\infty}a_j(\sigma)s^j,
\qquad m\in\Z,\qquad a_m(\sigma)\ne0.
\]
Its first nonzero Laurent coefficient is \(a_m(\sigma)\). The source writes \(\gamma^*(0,\sigma,\operatorname{Ad}_H)\) for a suitably normalized version of that coefficient, without specifying the normalization or the auxiliary choices defining the gamma factor. Thus \(\gamma^*\) cannot be identified unconditionally with \(a_m(\sigma)\) from this record alone.

\textbf{Supplied conjectural statement.} For every reductive group \(H\) in the intended, unspecified class, after fixing the missing parameter data and compatible measure and gamma-factor normalizations, prove that for almost every tempered representation \(\sigma\),
\[
 d_H(\sigma)
 =\frac{\dim(\rho_\sigma)}{|S_\sigma|}
   \left|\gamma^*(0,\sigma,\operatorname{Ad}_H)\right|.
\]
The exceptional set is required to have measure zero for the intended measure, which the record does not identify. These missing choices prevent this extracted formula from being a fully specified standalone conjecture.
\subsection{Short English statement}
The supplied conjecture predicts that, for almost every tempered representation, the Plancherel density equals an enhancement multiplicity factor times the absolute value of a normalized leading Laurent coefficient of an adjoint gamma factor. The group domain, parameter data, measures, and normalizations require clarification before the formula has a precise standalone meaning.
\subsection{Sources}
\begin{itemize}
\item Local registry: \texttt{data/unsolvedmath/problems.json}, with \texttt{catalog/problems.json} as the additional-record registry; supplied permanent UnsolvedMath ID \(30004774\), exact source problem code \texttt{catalogue entry}. These identifiers cite the supplied record; the local files were not independently inspected.
\item Supplied report citation: \emph{Automorphic Forms, Geometry and Arithmetic (2022)}, Oberwolfach Reports, pp.~2089--2156, DOI \texttt{10.4171/owr/2021/39}. \url{https://doi.org/10.4171/owr/2021/39}. The literature-source label is ``Oberwolfach Report 39/2021.'' The record gives proposed year 2021 and attributes the conjecture to Hiraga--Ichino--Ikeda--Langlands; no individual proposer is supplied.
\item Supplied literature source: \emph{Hiraga--Ichino--Ikeda conjecture for classical groups}. \url{https://arxiv.org/abs/2508.08470}. No authors or further bibliographic details are supplied in the record.
\item Literature assessment dated exactly \texttt{2026-08-22}: the supplied curation labels the problem partially solved and reports that the prediction is proved in the Archimedean case and in substantial \(p\)-adic families, while the extracted general Plancherel-density statement is not established for all reductive groups. This is supplied dated curation, not an independently verified literature assessment. The background expressly calls for independent verification of time-sensitive status and cited claims.
\item The supplied statement status is \texttt{corrected\_verified}, accompanied by a claim that the standalone statement and hypotheses were checked against the official report. That verification claim is part of the source record and has not been independently confirmed here; neither linked source was visited.
\end{itemize}
\end{document}
